% Generated draft: requires mathematical and source review.
\documentclass[11pt]{article}
\usepackage[margin=25mm]{geometry}
\usepackage{fontspec}
\setmainfont{Latin Modern Roman}
\usepackage{amsmath,amssymb,mathtools,unicode-math}
\setmathfont{Latin Modern Math}
\usepackage{enumitem,needspace,xurl,hyperref}
\setlength{\parindent}{0pt}
\setlength{\parskip}{5pt}
\setlength{\emergencystretch}{3em}
\newcommand{\N}{\mathbb{N}}
\newcommand{\Z}{\mathbb{Z}}
\newcommand{\Q}{\mathbb{Q}}
\newcommand{\R}{\mathbb{R}}
\newcommand{\C}{\mathbb{C}}
\newcommand{\F}{\mathbb{F}}
\newcommand{\A}{\mathbb{A}}
\newcommand{\PP}{\mathbb{P}}
\newcommand{\E}{\mathbb{E}}
\newcommand{\eps}{\varepsilon}
\newcommand{\dd}{\,\mathrm{d}}
\begin{document}
\section*{UnsolvedMath ID: 30005667}
\textbf{Source problem code:} catalogue entry
\subsection{Definitions and mathematical statement}
For a manifold \(M\) and a differentiability index \(t\), once the manifold structure, boundary conventions, and meaning of \(C^t\) have been specified, let \(\operatorname{Diff}^t(M)\) denote the group of bijections \(M\to M\) whose maps and inverses have regularity \(C^t\). An action of an abstract group \(\Gamma\) on \(M\) by \(C^t\) diffeomorphisms is a homomorphism \(\rho:\Gamma\to\operatorname{Diff}^t(M)\). Its kernel is \(\ker\rho=\{\gamma\in\Gamma:\rho(\gamma)=\operatorname{id}_M\}\). The action is trivial if \(\rho(\Gamma)=\{\operatorname{id}_M\}\), equivalently if \(\ker\rho=\Gamma\).

The supplied question asks whether there exists a countable subgroup \(\Gamma_{r,M}\leq\operatorname{Diff}^r(M)\), or, more strongly, a finitely generated such subgroup, for which every action on the same manifold \(M\) at every higher regularity \(s>r\) is trivial, or at least has a ``large kernel.'' The subgroup inclusion supplies a faithful action at regularity \(r\).

Conditional on specifying \(M\), \(r\), and the allowed differentiability indices, the stronger, trivial-action version is
\[
\exists\,\Gamma_{r,M}\leq\operatorname{Diff}^r(M),\quad
|\Gamma_{r,M}|\leq\aleph_0,\quad
\forall s>r\ \forall\rho\in\operatorname{Hom}\!\left(\Gamma_{r,M},\operatorname{Diff}^s(M)\right),\quad
\rho(\Gamma_{r,M})=\{\operatorname{id}_M\}.
\]
The finitely generated version additionally requires a finite set \(S\subseteq\Gamma_{r,M}\) with \(\Gamma_{r,M}=\langle S\rangle\). These are distinct requested strengths of the existence question, rather than restrictions on the particular faithful action supplied by the inclusion alone.

The source also permits a weaker conclusion that each higher-regularity action has a ``large kernel,'' but supplies no definition of largeness; this alternative cannot be given a precise kernel condition from the record. The range and quantification of \(M\) and \(r\), including dimension, compactness, boundary conditions, and nonintegral regularity conventions, are unspecified. The record also does not explicitly require \(\Gamma_{r,M}\neq\{1\}\); an intended nontriviality condition requires clarification.
\subsection{Short English statement}
The question seeks a countable, preferably finitely generated, group acting faithfully on \(M\) with regularity \(C^r\), while every action of that group on the same \(M\) at any higher regularity is trivial or has a suitably large kernel. The manifold hypotheses, regularity conventions, nontriviality requirement, and meaning of ``large'' remain unspecified.
\subsection{Sources}
\begin{itemize}
\item Supplied local UnsolvedMath registry record, permanent ID \(30005667\), exact source problem code \texttt{catalogue entry}, titled ``Group Actions with Differentiability Obstructions,'' in the Oberwolfach Reports Open Problems collection.
\item Supplied report citation: \emph{Teichm\"uller Theory: Classical, Higher, Super and Quantum} (2024), Oberwolfach Reports, pp.\ 1843--1892, DOI \texttt{10.4171/owr/2023/33}. \url{https://doi.org/10.4171/owr/2023/33}. The extracted passage identifies p.\ 1869 and the contribution ``First order theory of homeomorphism groups of manifolds, and applications,'' by Thomas Koberda, joint work with Sang-hyun Kim and J.\ de la Nuez Gonz\'alez. The registry records proposed year 2023 and supplies no separate proposer attribution.
\item Kim--Koberda, \emph{Structure and regularity of group actions}. Supplied literature link: \url{https://arxiv.org/abs/2106.15532}.
\item Koberda, \emph{Diffeomorphism groups of critical regularity}. Supplied literature link: \url{https://doi.org/10.1090/noti2252}.
\item Supplied literature assessment dated \texttt{2026-08-22}: the record labels the problem ``partially solved'' and states that the proposed regularity obstruction is proved for compact one-manifolds, including finitely generated examples, while the analogous higher-dimensional existence question remains open. This is supplied dated curation derived from the OWR record, not an independently verified literature assessment; the precise hypotheses and strength of the cited results require checking.
\item The registry labels the statement \texttt{corrected\_verified} and reports that the official report was checked directly. That verification claim is supplied provenance and has not been independently verified in this draft.
\end{itemize}
\end{document}
